% Generated draft: requires mathematical and source review.
\documentclass[11pt]{article}
\usepackage[margin=25mm]{geometry}
\usepackage{fontspec}
\setmainfont{Latin Modern Roman}
\usepackage{amsmath,amssymb,mathtools,unicode-math}
\setmathfont{Latin Modern Math}
\usepackage{enumitem,needspace,xurl,hyperref}
\setlength{\parindent}{0pt}
\setlength{\parskip}{5pt}
\setlength{\emergencystretch}{3em}
\newcommand{\N}{\mathbb{N}}
\newcommand{\Z}{\mathbb{Z}}
\newcommand{\Q}{\mathbb{Q}}
\newcommand{\R}{\mathbb{R}}
\newcommand{\C}{\mathbb{C}}
\newcommand{\F}{\mathbb{F}}
\newcommand{\A}{\mathbb{A}}
\newcommand{\PP}{\mathbb{P}}
\newcommand{\E}{\mathbb{E}}
\newcommand{\eps}{\varepsilon}
\newcommand{\dd}{\,\mathrm{d}}
\begin{document}
\section*{UnsolvedMath ID: 801088}
\textbf{Source problem code:} catalogue entry
\subsection{Definitions and mathematical statement}
Let \(d\geq 1\) be an integer, and let \(f\colon\C^d\to\C^d\) be a polynomial automorphism, meaning a polynomial map with a polynomial inverse. For a polynomial map \(g=(g_1,\ldots,g_d)\), let \(\deg(g)=\max_i\deg(g_i)\), where each coordinate degree is total degree. For \(n\geq 1\), write \(f^n\) for the \(n\)-fold iterate of \(f\). Interpret the notation \(\lambda(f)\) in the supplied question as the first dynamical degree
\[
\lambda(f)=\lim_{n\to\infty}\bigl(\deg(f^n)\bigr)^{1/n};
\]
this limit exists by submultiplicativity of the degrees of iterates.

Suppose that
\[
\lambda(f)=1,
\qquad
\sup_{n\geq 1}\deg(f^n)=\infty.
\]
Must there exist a constant \(C=C(f)>0\) such that
\[
\deg(f^n)\geq Cn\qquad\text{for every integer }n\geq 1?
\]
The dimension is unrestricted in the supplied statement; the question is therefore formulated for arbitrary finite \(d\geq 1\). The constant may depend on the automorphism, including its dimension. The definition of \(\lambda(f)\) above makes explicit the conventional interpretation of notation left undefined in the supplied excerpt.

For context, the supplied source attributes to Urech, under the reference label \([\mathrm{urech}]\), the assertion that unbounded degree growth tends to infinity. It attributes a stronger assertion to Cantat and Xie, under \([\mathrm{CX}]\): there exists a universal function \(\sigma\colon\N\to\N\) such that
\[
\limsup_{n\to\infty}\sigma(n)=\infty,
\qquad
\deg(f^n)\geq\sigma(n).
\]
These are reported background assertions, not independently verified here. The excerpt does not specify the full hypotheses or scope of universality of this background bound. It describes the present question as a weaker form of a preceding problem, whose statement is not supplied.
\subsection{Short English statement}
For a polynomial automorphism of \(\C^d\) with first dynamical degree \(1\) and unbounded degrees of iterates, must those degrees have a positive linear lower bound for every positive iterate? The constant may depend on the automorphism. The supplied registry labels this question open; this draft does not independently verify that status.
\subsection{Sources}
\begin{itemize}
\item Local UnsolvedMath registry, permanent ID \(801088\), exact source problem code \texttt{catalogue entry}; supplied complete source record, upstream revision \texttt{d371ee917ac9647fbb06cd6fc529deec6f3dda82}. The registry attributes the problem to Berger and Bedford, with proposed year 2023, and records its status as open.
\item Berger and Bedford, \emph{Henon Maps: A List of Open Problems} (2023), Question 40. Supplied source URL: \url{https://arxiv.org/abs/2312.03907}. The record reports access on 2026-10-16 and extraction by \texttt{arxiv\_latex}; these are supplied provenance details, not a claim of independent consultation or verification.
\item The supplied excerpt attributes background degree-growth results to Urech, reference label \([\mathrm{urech}]\), and to Cantat and Xie, reference label \([\mathrm{CX}]\). Full bibliographic entries and separate URLs for these references are not supplied.
\end{itemize}
\end{document}
